%%%%%%%%%%%%%%%%%%%%%%%%%%%%%%%%%%%%%%%%%%%%%%%%%%%%%%%%%%%%%%%%%%%%%%
% Lớp 9 | Chương X. Một số hình khối trong thực tiễn | Hình trụ và hình nón
% Dạng 9H0?31-3: Bài toán thực tế về hình trụ, hình nón
%%%%%%%%%%%%%%%%%%%%%%%%%%%%%%%%%%%%%%%%%%%%%%%%%%%%%%%%%%%%%%%%%%%%%%


% nguồn: 2026-11-14_09_de_tuyen_sinh_10_tphcm_2025_2026 (nhập 14/11/2026)
\begin{bt}%[9H0V31-3]%[ma:7f59a503]
% nguon: Đề thi chính thức Tuyển sinh 10 TP.HCM năm học 2025-2026 – Bài 5
Một hộp đựng bóng tennis có dạng hình trụ chứa vừa khít $4$ quả bóng tennis có dạng hình cầu. Biết diện tích bề mặt mỗi quả bóng tennis là $132{,}67\text{ cm}^2$.
Cho biết diện tích bề mặt hình cầu bán kính $R$ là $S = 4\pi R^2$. Diện tích toàn phần hình hộp chữ nhật có ba kích thước $a, b, c$ là $S_{tp} = 2(ab + bc + ca)$.
\begin{enumerate}[a)]
	\item Tính bán kính của mỗi quả bóng tennis.
	\item Nhà sản xuất sử dụng các thùng giấy hình hộp chữ nhật (có nắp) để chứa $12$ hộp tennis sao cho các hộp tennis được xếp vừa khít trong thùng giấy (xếp thành $3$ hàng, mỗi hàng $4$ hộp; hoặc $2$ hàng, mỗi hàng $6$ hộp theo chiều đứng). Hỏi cần tối thiểu bao nhiêu $\text{m}^2$ giấy để thiết kế một thùng như trên (giả sử mép dán không đáng kể). Các kết quả được làm tròn đến hàng phần trăm.
\end{enumerate}
\loigiai{
\begin{enumerate}[a)]
	\item Diện tích bề mặt mỗi quả bóng tennis là $132{,}67\text{ cm}^2$, do đó:
	$4\pi R^2 = 132{,}67 \Rightarrow R^2 = \dfrac{132{,}67}{4\pi} \approx \dfrac{132{,}67}{4\cdot 3{,}14159} \approx 10{,}557$.\\
	Suy ra bán kính của mỗi quả bóng tennis là:
	$R = \sqrt{\dfrac{132{,}67}{4\pi}} \approx 3{,}25\text{ (cm)}$.
	\item Đường kính của mỗi quả bóng tennis là:
	$d = 2R = 2\cdot 3{,}25 = 6{,}5\text{ (cm)}$.\\
	Vì hộp tennis hình trụ chứa vừa khít $4$ quả bóng nên:
	\begin{itemize}
		\item Đường kính đáy hộp hình trụ là: $D = d = 6{,}5\text{ (cm)}$.
		\item Chiều cao của hộp hình trụ là: $h = 4d = 4\cdot 6{,}5 = 26\text{ (cm)}$.
	\end{itemize}
	Để chứa vừa khít $12$ hộp tennis dựng đứng thành khối hình hộp chữ nhật, cách xếp tối ưu diện tích là xếp thành $3$ hàng, mỗi hàng $4$ hộp (hoặc đáy gồm $3 \times 4$ hộp):
	\begin{itemize}
		\item Chiều rộng thùng giấy là: $a = 3\cdot D = 3\cdot 6{,}5 = 19{,}5\text{ (cm)}$.
		\item Chiều dài thùng giấy là: $b = 4\cdot D = 4\cdot 6{,}5 = 26\text{ (cm)}$.
		\item Chiều cao thùng giấy là: $c = h = 26\text{ (cm)}$.
	\end{itemize}
	Diện tích toàn phần của thùng giấy hình hộp chữ nhật là:
	$S_{tp} = 2(ab + bc + ca) = 2(19{,}5\cdot 26 + 26\cdot 26 + 26\cdot 19{,}5) = 2(507 + 676 + 507) = 2\cdot 1690 = 3380\text{ (cm}^2\text{)}$.\\
	Đổi sang đơn vị mét vuông:
	$3380\text{ cm}^2 = \dfrac{3380}{10000}\text{ m}^2 = 0{,}338\text{ m}^2 \approx 0{,}34\text{ (m}^2\text{)}$.\\
	Vậy cần tối thiểu khoảng $0{,}34\text{ m}^2$ giấy để thiết kế một thùng như trên.
\end{enumerate}
}
\end{bt}
